\documentclass[11pt,a4paper]{article}
\usepackage[T1]{fontenc}
\usepackage[utf8]{inputenc}
\usepackage[main=italian,provide=*]{babel}
\usepackage{amsmath,amssymb}
\usepackage{geometry}
\geometry{margin=2.3cm,top=2.5cm,bottom=2.7cm}
\usepackage{booktabs}
\usepackage[table]{xcolor}
\usepackage{tcolorbox}
\tcbuselibrary{skins,breakable}
\usepackage{titlesec}
\usepackage{enumitem}
\usepackage{needspace}
\usepackage{fancyhdr}
\usepackage{hyperref}
\hypersetup{colorlinks=true,linkcolor=Sepia,citecolor=Sepia,urlcolor=Sepia}

% --- palette ---
\definecolor{inkblue}{HTML}{1B3A5C}
\definecolor{lightblue}{HTML}{EAF1F8}
\definecolor{accent}{HTML}{B0562A}
\definecolor{softgray}{HTML}{6B6B6B}
\definecolor{okgreen}{HTML}{2E6B3E}
\definecolor{okgreenbg}{HTML}{EAF5EC}
\definecolor{warnamber}{HTML}{9C6B12}
\definecolor{warnbg}{HTML}{FBF1E0}

% --- titoli di sezione ---
\titleformat{\section}
  {\normalfont\Large\bfseries\color{inkblue}}
  {\thesection}{0.6em}{}[{\color{inkblue!40}\titlerule}]
\titlespacing{\section}{0pt}{0.7em}{0.4em}
\titleformat{\subsection}
  {\normalfont\large\bfseries\color{accent}}
  {}{0em}{}
\titlespacing{\subsection}{0pt}{0.5em}{0.2em}

\pagestyle{fancy}
\fancyhf{}
\renewcommand{\headrulewidth}{0.4pt}
\fancyhead[L]{\small\color{softgray}\texttt{vqe\_dimer.py} --- spiegazione e risultati}
\fancyhead[R]{\small\color{softgray}\thepage}
\fancyfoot[C]{\small\color{softgray}Tesi triennale, Samuele Schiavi --- Relatore: Prof.\ Alessandro Chiesa}

% --- box riutilizzabili ---
\newtcolorbox{keyeq}[1][]{
  colback=lightblue, colframe=inkblue!70, boxrule=0.6pt,
  arc=2pt, left=8pt, right=8pt, top=3pt, bottom=3pt, #1
}
\newtcolorbox{panelbox}[1]{
  colback=white, colframe=softgray!50, boxrule=0.5pt,
  arc=2pt, left=7pt, right=7pt, top=2pt, bottom=2pt,
  fonttitle=\bfseries\color{inkblue}, title={#1}
}
\newtcolorbox{okbox}[1]{
  colback=okgreenbg, colframe=okgreen!60, boxrule=0.6pt,
  arc=2pt, left=7pt, right=7pt, top=2pt, bottom=2pt,
  fonttitle=\bfseries\color{okgreen}, title={#1}
}
\newtcolorbox{warnbox}[1]{
  colback=warnbg, colframe=warnamber!70, boxrule=0.6pt,
  arc=2pt, left=7pt, right=7pt, top=2pt, bottom=2pt,
  fonttitle=\bfseries\color{warnamber}, title={#1}
}

\title{\vspace{-1.5em}\color{inkblue}\texttt{vqe\_dimer.py}\\[0.2em]
\Large\normalfont\color{softgray}Spiegazione concettuale e lettura dei risultati}
\author{Note di lavoro --- Tesi triennale, Samuele Schiavi\\ Relatore: Prof.\ Alessandro Chiesa}
\date{}

\begin{document}
\sloppy
\maketitle
\thispagestyle{fancy}
\vspace{-2.6em}

Il modulo implementa il VQE per il dimero e poggia sul benchmark esatto
(\texttt{dimer\_exact.py}), che fornisce il riferimento contro cui
misurare l'errore.

\section*{Principio variazionale (perché il VQE funziona)}
\addcontentsline{toc}{section}{Principio variazionale}

Per qualsiasi stato normalizzato $|\psi\rangle$ vale
$\langle\psi|H|\psi\rangle \geq E_0$, con uguaglianza solo se
$|\psi\rangle$ è il ground state. Parametrizzando una famiglia
$|\psi(\boldsymbol\theta)\rangle$ e \textbf{minimizzando}
$\langle H\rangle_{\boldsymbol\theta}$, il VQE si avvicina a $E_0$
\textbf{da sopra}: non può mai scendere sotto. È il motivo per cui
$\Delta E = E_\text{VQE} - E_0 \geq 0$ sempre (a meno del rumore
numerico). Nel notebook questo si verifica numericamente: stati casuali
hanno tutti $\langle H\rangle \geq E_0$.

\subsection*{Estimator (come si calcola $\langle H\rangle$)}

A ogni iterazione del ciclo quantistico-classico,
\texttt{StatevectorEstimator} valuta
$\langle\psi(\boldsymbol\theta)|H|\psi(\boldsymbol\theta)\rangle$. Si
costruisce un \textbf{PUB} \texttt{(ansatz, hamiltonian, [params])}; il
risultato sta in \texttt{result[0].data.evs}, un array NumPy --- serve
\texttt{.item()} per ottenere uno scalare. Il notebook mostra che
l'estimator dà lo stesso valore del calcolo manuale via
\texttt{Statevector}.

\section*{I due ansätze}
\addcontentsline{toc}{section}{I due ansätze}

\begin{panelbox}{HA --- hardware-efficient (\texttt{n\_local}, $R_y$+CZ)}
Generico: non sa nulla della fisica del problema, ma con abbastanza
strati rappresenta qualunque stato. Parametri: $2(\text{reps}+1)$, cioè
\textbf{6} per \texttt{reps=2}. (\texttt{n\_local} sostituisce la classe
\texttt{TwoLocal}, deprecata in Qiskit 2.1.)
\end{panelbox}
\smallskip

\begin{panelbox}{PMA --- physically motivated (Crippa et al.\ 2021)}
Costruito \textbf{sulle simmetrie} del problema. Blocco fondamentale su
ogni legame:
$$W_{ij}(\theta) = \text{CNOT}_{ij}\,\big(R_y(\theta)_i\otimes I_j\big)\,\text{CNOT}_{ij}.$$
Lo stato iniziale rispetta la simmetria del ground state atteso:
\begin{itemize}[leftmargin=1.3em,itemsep=0.1em]
\item $B/J<2$ (GS = singoletto, $S=0,M=0$): si parte dal singoletto
locale $(|01\rangle-|10\rangle)/\sqrt2$;
\item $B/J\geq2$ (GS = $|11\rangle$, $M=-1$): si parte da $|11\rangle$.
\end{itemize}
Per il dimero ($N=2$, un legame, $K=1$): \textbf{1 solo parametro}. Il
prezzo dell'efficienza è che il PMA \textbf{conserva $S$ e $M$}.
\end{panelbox}

\subsection*{Perché si sceglie il riferimento in base a $B/J$}

Il ground state del dimero \textbf{cambia natura bruscamente} al livello
critico $B/J=2$ (level crossing): sotto soglia è il singoletto, sopra è
$|11\rangle$ --- due stati completamente diversi, non c'è un singolo
riferimento che vada bene ovunque. L'ansatz quindi \textbf{sceglie il
ramo} in base al regime, invece di usare un unico circuito che copra
continuamente tutto l'asse $B/J$ (come si fa in
\texttt{confronto\_ansatz\_entangler.ipynb} con il blocco RBS/Givens,
"sicuro" anche partendo da una sovrapposizione dei due settori). Le due
scelte sono entrambe legittime; qui la si vede esplicitamente passo per
passo.

\subsection*{Derivazione esplicita: perché $X,H,\mathrm{CX},X$ prepara il singoletto}

Seguendo lo stato passo per passo, partendo da $|00\rangle$ (convenzione
$|q_0\,q_1\rangle$):

\begin{center}
\renewcommand{\arraystretch}{1.25}
\begin{tabular}{@{}ll@{}}
\toprule
\rowcolor{inkblue!10}
\textbf{Gate} & \textbf{Stato dopo}\\
\midrule
(iniziale) & $|00\rangle$\\
$X(q_0)$ & $|10\rangle$\\
$H(q_0)$ & $\frac{1}{\sqrt2}\big(|00\rangle-|10\rangle\big)$\\
$\mathrm{CX}(q_0\to q_1)$ & $\frac{1}{\sqrt2}\big(|00\rangle-|11\rangle\big)$ (stato di Bell)\\
$X(q_0)$ & $\frac{1}{\sqrt2}\big(|10\rangle-|01\rangle\big)=-\frac{1}{\sqrt2}\big(|01\rangle-|10\rangle\big)$\\
\bottomrule
\end{tabular}
\end{center}

L'ultimo stato è il \textbf{singoletto}, a meno di una fase globale $-1$
irrilevante. Quattro gate fissi, senza parametri: non è un ansatz
variazionale, è una \textbf{preparazione deterministica}. Il notebook
verifica ogni stadio intermedio con il simulatore (non solo il risultato
finale), così si vede esattamente dove nasce l'entanglement (al passo
CX, che produce prima uno stato di Bell) e dove arriva il singoletto
(all'ultimo $X$).

\subsection*{Perché basta 1 parametro: la struttura a blocchi di $W_{01}(\theta)$}

Il blocco $W_{01}(\theta)$ \textbf{non} è semplicemente "l'identità
fuori dal settore $M=0$", come si potrebbe pensare. È qualcosa di più
preciso, verificato nel notebook stampando la sua matrice $4\times4$
esplicita: è \textbf{diagonale a blocchi}, con la \textbf{stessa}
rotazione (stesso angolo $\theta$) che agisce \textbf{sia} su
$\{|01\rangle,|10\rangle\}$ \textbf{sia} su $\{|00\rangle,|11\rangle\}$,
senza mai farli comunicare:

\begin{keyeq}
\centering
$W_{01}(\theta) = \begin{pmatrix} R(\theta) & 0 \\ 0 & R(\theta)\end{pmatrix}$
\ nei due sottospazi $\{|00\rangle,|11\rangle\}$ e $\{|01\rangle,|10\rangle\}$.
\end{keyeq}

\textbf{Perché questo spiega tutto.} Il blocco non fa uscire lo stato
dal sottospazio invariante in cui lo ha messo la preparazione iniziale:
\begin{itemize}[leftmargin=1.3em,itemsep=0.1em]
\item se parti dal \textbf{singoletto} (ampiezza solo su
$\{|01\rangle,|10\rangle\}$), il blocco fa scorrere lo stato dentro
\emph{quel} sottospazio (dal singoletto verso il tripletto $M=0$ e
viceversa);
\item se parti da $|11\rangle$ (ampiezza solo su
$\{|00\rangle,|11\rangle\}$), il blocco fa scorrere lo stato dentro
\emph{quell'altro} sottospazio.
\end{itemize}

I due sottospazi non si mescolano mai (la matrice è a blocchi), ma il
blocco \textbf{non} è l'identità su nessuno dei due --- li ruota
entrambi, con lo stesso angolo. Ecco perché \textbf{un solo parametro
$\theta$} basta per esplorare correttamente entrambi i rami: non
servono due parametri separati, perché la stessa rotazione
riparametrizza correttamente qualunque dei due sottospazi ti serva, una
volta scelto il ramo dalla preparazione.

\textbf{Confronto con la rotazione di Givens/RBS "vera".} Nel notebook
\texttt{confronto\_ansatz\_entangler.ipynb} si usa invece un blocco che
è \textbf{l'identità} su $\{|00\rangle,|11\rangle\}$ (li lascia fermi
del tutto), non un'altra rotazione con lo stesso angolo. Sono due
decomposizioni diverse della stessa idea --- "un solo parametro per il
settore $M=0$" --- con proprietà leggermente diverse: quella di questo
notebook richiede di scegliere il ramo giusto in partenza (perché
ruoterebbe anche l'altro sottospazio se ci fosse ampiezza lì); quella di
Givens/RBS sarebbe "sicura" anche partendo da una sovrapposizione dei
due settori. Nessuna delle due è più corretta: sono scelte di
implementazione diverse dello stesso principio fisico, ed entrambe
compaiono nel progetto per un confronto esplicito.

\subsection*{Fidelity}

L'energia dice quanto sei vicino in \emph{valore}; la fidelity dice
quanto sei vicino in \emph{stato}:
$$\mathcal F = |\langle\psi_0|\psi(\tilde{\boldsymbol\theta})\rangle|,$$
con $|\psi_0\rangle$ il ground state esatto da \texttt{eigh}.
$\mathcal F=1$ significa che il VQE ha trovato esattamente il ground
state, non solo la sua energia.

\needspace{4\baselineskip}
\section*{Configurazione della run}
\addcontentsline{toc}{section}{Configurazione della run}

Confronto incrociato \textbf{2 ansätze $\times$ 2 valori di $D$}, su
griglia $4\times4$ pannelli (righe = configurazioni, colonne =
metriche). Output salvato in \texttt{.png} (raster) e \texttt{.pdf}
(vettoriale, zoomabile).

\begin{center}
\renewcommand{\arraystretch}{1.25}
\begin{tabular}{@{}cp{5.3cm}cc@{}}
\toprule
\rowcolor{inkblue!10}
\textbf{Riga} & \textbf{Ansatz} & \textbf{$D$} & \textbf{Parametri}\\
\midrule
1 & HA (hardware-efficient, \texttt{n\_local} $R_y$+CZ, reps=2) & 0 & 6\\
2 & HA & 0.2 & 6\\
3 & PMA (physically motivated, Crippa et al., $K=1$) & 0 & 1\\
4 & PMA & 0.2 & 1\\
\bottomrule
\end{tabular}
\end{center}

Parametri comuni: $J=1$, 20 punti $B/J\in[0,5]$, ottimizzatore
L-BFGS-B, 3 restart.

\begin{keyeq}
\small
Nota: il PMA per il dimero con $K=1$ ha \textbf{1 parametro} (un solo
legame, un solo gate $W_{01}(\theta)$), non 2. Lo stato iniziale dipende
da $B/J$: singoletto locale sotto $B/J=2$, $|11\rangle$ sopra.
\end{keyeq}

\needspace{4\baselineskip}
\section*{Risultato riga per riga}
\addcontentsline{toc}{section}{Risultato riga per riga}

\begin{okbox}{Riga 1 --- HA, $D=0$ (baseline)}
\textbf{Energia:} VQE sovrapposto all'esatto su tutto il range.\ \
\textbf{Errore:} $\Delta E\sim10^{-13}$--$10^{-14}$, rumore numerico
puro.\ \ \textbf{Fidelity:} $\mathcal F=1.000$ ovunque.\ \
\textbf{Magnetizzazione:} salto netto $0\to-1$ in $B/J=2$, riprodotto
esattamente. L'HA con 6 parametri è espressivamente completo per il
dimero isotropo.
\end{okbox}
\smallskip

\begin{okbox}{Riga 2 --- HA, $D=0.2$}
\textbf{Energia e fidelity:} identiche qualità del caso $D=0$ ---
$\Delta E\sim10^{-13}$, $\mathcal F=1.000$ ovunque, \textbf{anche
attraverso l'anticrossing}.\ \ \textbf{Magnetizzazione:} crossover
liscio attraverso $B/J=2$ (non più un gradino), VQE sovrapposto
all'esatto. L'HA, essendo generico, rappresenta senza problemi il
ground state misto dell'anticrossing: il termine DM non lo penalizza.
\end{okbox}
\smallskip

\begin{okbox}{Riga 3 --- PMA, $D=0$}
\textbf{Errore:} $\Delta E\sim10^{-16}$ --- precisione macchina,
\textbf{migliore dell'HA}.\ \ \textbf{Fidelity:} $\mathcal F=1.000$
ovunque.\ \ \textbf{Magnetizzazione:} salto netto perfetto. Il PMA
raggiunge il ground state esatto con \textbf{1 solo parametro} contro i
6 dell'HA: è il vantaggio centrale di Crippa et al., imporre le
simmetrie del problema ($S,M$ conservati) riduce drasticamente la
complessità variazionale.
\end{okbox}
\smallskip

\begin{warnbox}{Riga 4 --- PMA, $D=0.2$ (il caso istruttivo)}
\textbf{Errore:} sale fino a $\Delta E\simeq0.2$ in $B/J=2$, con un
picco netto.\ \ \textbf{Fidelity:} scende a $\mathcal F\simeq0.82$ in
$B/J=2$, poi risale.\ \ \textbf{Magnetizzazione:} resta un
\textbf{gradino netto} $0\to-1$, NON il crossover liscio dell'esatto.

Questo \textbf{non è un bug}: è la conseguenza diretta della struttura
del PMA. Il PMA costruisce per costruzione stati a $S,M$ definiti
(singoletto $M=0$ sotto soglia, $|11\rangle$ con $M=-1$ sopra). Ma con
$D\neq0$ il vero ground state \textbf{non ha più $M$ definito} --- è la
miscela singoletto/tripletto $M=-1$ dell'anticrossing. Il PMA, che
conserva $M$, non può rappresentare quella miscela. La degradazione è
massima in $B/J=2$, dove il mescolamento dei due settori è massimo, e
svanisce lontano dalla soglia. La magnetizzazione resta un gradino
perché gli stati del PMA hanno $\langle M_z\rangle$ esattamente $0$ o
$-1$, mai valori intermedi.
\end{warnbox}

\needspace{4\baselineskip}
\section*{Il risultato centrale}
\addcontentsline{toc}{section}{Il risultato centrale}

\begin{center}
\renewcommand{\arraystretch}{1.3}
\begin{tabular}{@{}lccc@{}}
\toprule
\rowcolor{inkblue!10}
\textbf{Configurazione} & \textbf{Parametri} & \textbf{Accuratezza} & \textbf{Fidelity}\\
\midrule
HA, $D=0$   & 6 & $\Delta E\sim10^{-13}$ & $\mathcal F=1$ ovunque\\
HA, $D=0.2$ & 6 & $\Delta E\sim10^{-13}$ & $\mathcal F=1$ ovunque\\
PMA, $D=0$  & 1 & $\Delta E\sim10^{-16}$ & $\mathcal F=1$ ovunque\\
PMA, $D=0.2$& 1 & $\Delta E$ fino a $\sim0.2$ in $B/J{=}2$ & $\mathcal F\simeq0.82$ in $B/J{=}2$\\
\bottomrule
\end{tabular}
\end{center}

Con $D=0$ il PMA \textbf{batte} l'HA: raggiunge il ground state esatto
(precisione macchina) con \textbf{1 parametro invece di 6}. È il
vantaggio centrale di Crippa et al.: imporre le simmetrie del problema
riduce drasticamente la complessità variazionale.

Con $D=0.2$ il PMA \textbf{crolla} vicino a $B/J=2$, perché il vero
ground state è la miscela
$$|\psi_0\rangle = \alpha\,|\text{singoletto}\rangle + \beta\,|11\rangle,$$
che il PMA --- conservando $M$ --- non può rappresentare. L'HA,
generico, non ha questo vincolo e resta a $\mathcal F=1$ in entrambi i
casi.

\begin{keyeq}
\textbf{Nota di collegamento.} Il notebook
\texttt{confronto\_ansatz\_entangler.ipynb} dimostra \emph{perché}
esattamente questo accade, per contrasto: mettendo accanto quattro
ansatz reali generici che non degradano affatto a $D\neq0$, isola la
rottura di simmetria $M$ come unica causa (non la complessità dello
stato), e mostra come ripararla --- con due famiglie di PMA estesi
(PMA-1q, PMA-2q) che restano nel campo reale e recuperano
$\mathcal F=1$ ovunque.
\end{keyeq}

\subsection*{Sintesi del confronto HA vs PMA}

\begin{center}
\renewcommand{\arraystretch}{1.3}
\begin{tabular}{@{}lcc@{}}
\toprule
\rowcolor{inkblue!10}
 & \textbf{HA} & \textbf{PMA}\\
\midrule
parametri (dimero) & 6 & 1\\
accuratezza $D=0$ & $\sim10^{-13}$ & $\sim10^{-16}$\\
accuratezza $D=0.2$ & $\sim10^{-13}$ (perfetta) & degrada a $B/J{=}2$ ($\mathcal F\simeq0.82$)\\
robustezza a rottura di simmetria & alta (generico) & bassa (vincolato a $S,M$)\\
\bottomrule
\end{tabular}
\end{center}

\textbf{Il trade-off è il messaggio centrale:} il PMA è enormemente più
efficiente (1 parametro vs 6, precisione macchina) \textbf{finché il
modello rispetta le simmetrie su cui è costruito}. Quando un termine
come il DM rompe la conservazione di $M$, il PMA fallisce proprio dove
la rottura è massima, mentre l'HA generico resta accurato. Questo è
esattamente il punto fisico che lega:
\begin{itemize}[leftmargin=1.3em,itemsep=0.1em]
\item la derivazione del termine DM (\texttt{dimero\_sintesi.pdf}),
\item l'analisi delle simmetrie ($[H_D,S_z]\neq0$),
\item e la scelta dell'ansatz.
\end{itemize}

\section*{Nota tecnica: warning "default.json not found"}
\addcontentsline{toc}{section}{Nota tecnica}

\texttt{show\_circuit} disegna i circuiti con
\texttt{qc.draw("mpl", style="iqp")}, uno stile esplicito esistente. Le
versioni recenti di Qiskit non includono più il file
\texttt{default.json} che il drawer cercherebbe con lo stile implicito,
generando altrimenti un \texttt{UserWarning} innocuo ma fastidioso ad
ogni disegno. Passare uno stile esistente (\texttt{iqp},
\texttt{textbook}, \texttt{clifford}) evita il problema alla radice.

\needspace{4\baselineskip}
\section*{La lezione e il ruolo nel progetto}
\addcontentsline{toc}{section}{La lezione}

Le simmetrie nell'ansatz sono \textbf{un'arma a doppio taglio}: il PMA è
enormemente più efficiente (1 vs 6 parametri, precisione macchina)
\textbf{finché il modello rispetta le simmetrie su cui è costruito};
quando un termine come il DM le rompe, il PMA fallisce proprio dove la
rottura è massima. Questo lega tre cose costruite separatamente: la
derivazione del termine DM, l'analisi delle simmetrie
($[H_D,S_z]\neq0$), e la scelta dell'ansatz --- inclusa la struttura
interna del blocco $W_{ij}(\theta)$, che rivela \emph{come} un solo
parametro riesca a coprire entrambi i rami del level crossing.

È un risultato completo e autocontenuto per il dimero $N=2$: stabilisce
il metodo di confronto (HA vs PMA, $D=0$ vs $D=0.2$) che verrà
replicato su $N=3$ nelle due topologie. Sarà particolarmente rilevante
per il \textbf{triangolo frustrato}, dove la struttura fisica
dell'ansatz conta di più.

\end{document}
